% META: {"id":"1SPE-VARALEA-CO-049","chapitre":"1SPE-VARIABLES-ALEATOIRES","type_objet":"corrige","exercice_ref":"1SPE-VARALEA-EX-049","capacites_codes":[],"extension_codes":["X1"],"programme_alignment":"OPTIONAL_EXTENSION","extension_label":"Approfondissement — Vers la Terminale","sources_inspiration":[],"status":"approved","origin":"generated","status_history":["generated","audited","accepted","approved"],"release_acceptance":"RELEASE_OWNER_BATCH_ACCEPTANCE","acceptance_closure_digest":"sha256:9b3ccf9a81c5520fbb7e03b7b2d3e7bf2057a02834b6d908bf3be5f49f3b553f","acceptance_receipt_digest":"sha256:4fad86f0282e0fa4e0419cca1ad6379ae7af5a280c04a4a90880f90a1c044242"}
\begin{center}\textbf{Approfondissement — Vers la Terminale}\end{center}
% BEGIN-VERIFY
% from sympy import Rational
% p = Rational(9,10); n = 8
% P8 = p**8
% E = n*p; assert E == Rational(36,5)
% END-VERIFY
\begin{corrige}{1SPE-VARALEA-EX-049}

\textbf{1.} Chaque vaccination est une épreuve de Bernoulli ($p = 0{,}9$), les $8$ sont indépendantes : $X \sim \mathcal{B}(8, 0{,}9)$.

\medskip
\textbf{2.} $P(X=8) = (0{,}9)^8 \approx 0{,}430$.

\medskip
\textbf{3.} $E(X) = 8 \times 0{,}9 = 7{,}2$ personnes protégées en moyenne.

\end{corrige}
